\documentclass{article}
\usepackage[utf8]{inputenc}
\usepackage[english, spanish]{babel}
%\usepackage[dvips]{graphics}
\usepackage{amsmath}
\usepackage{amssymb}
\usepackage{fullpage}
\usepackage{epsfig}
\usepackage{multicol}
\usepackage{wasysym}
\usepackage{tikz}
\usepackage{graphicx}
\usepackage{mathtools}
\usepackage{enumitem}
\usepackage[mathscr]{euscript}

\usetikzlibrary{arrows,decorations.pathmorphing,backgrounds,positioning,fit,calc,automata}
\usetikzlibrary{decorations.pathreplacing}
\usetikzlibrary{trees}
\usetikzlibrary{shapes}
\usetikzlibrary{chains}
\usetikzlibrary{patterns}
\usetikzlibrary{matrix}
%\usepackage{stackengine}
\pagenumbering{gobble}

%\usepackage[usenames,dvipsnames]{xcolor}
\usepackage{hyperref}
\definecolor{linkcolour}{rgb}{0,0.2,0.6}
\hypersetup{colorlinks,breaklinks,urlcolor=linkcolour,linkcolor=linkcolour}

\parindent 0pt

\newcommand{\cL}{\mathcal{L}}


\begin{document}

\includegraphics[width=1.8cm]{uc.png}
\vspace*{-1.9cm}

\hspace*{2.1cm}
 \begin{tabular}{l}
  \sc Pontificia Universidad Católica de Chile \\
  \sc Escuela de Ingeniería \\
  \sc Departamento de Ciencia de la Computación
  %\vspace{15\baselineskip}\mbox{}
  %\vspace{-3mm}\mbox{}
 \end{tabular}
 \bigskip

\vspace*{5mm}
\begin{center}
{IIC2213 --- Lógica para ciencia de la computación --- 1' 2023} \\
\vspace{3mm}
{\Large\bf TAREA 6} \\
\vspace{2mm}
\end{center}

\begin{tabular}{ll}
Publicación: & Martes 6 de junio. \\
Entrega: & \textbf{Lunes 19 de junio hasta las 23:59 horas}. \\
\end{tabular}

\subsection*{Indicaciones}

\begin{itemize}
\item Cada pregunta tiene 6 puntos ($+1$ base) y la nota de la tarea es el promedio de  las preguntas.
\item La solución debe estar escrita en \LaTeX. No se aceptarán tareas escritas de otra forma.
\item La tarea es individual, pudiendo discutirla con sus pares. Toda referencia externa debe citarse.
\end{itemize}

\subsection*{Objetivos}
\begin{itemize}
    \item Aplicar LPO al contexto de consultas en bases de datos.
    \item Aplicar conceptos de teorías.
\end{itemize}

\subsection*{Pregunta 1: Aplicación de LPO a consultas}
Decimos que dos grafos $G_1 = (V_1, E_1)$ y $G_2 =
(V_2, E_2)$\ son {\em isomorfos}, denotado por $G_1\cong G_2$, si existe una biyección $f : V_1
\rightarrow V_2$ tal que para todo $a,b\in V_1$ se tiene que
$$\{a,b\} \in E_1\quad\text{si y solo si}\quad\{f(a), f(b)\} \in E_2$$ 

Dado un par de grafos no dirigidos $G_1=(V_1,E_1)$ y $G_2=(V_2,E_2)$, construya una fórmula proposicional $\varphi_\text{iso}$ tal que 
$$\varphi_\text{iso}\text{ es satisfacible\quad si y solo si}\quad G_1\cong G_2$$
Demuestre la correctitud de su construcción.
\subsubsection*{Solución P1.}
Aquí va mi solución
\newpage

\subsection*{Pregunta 2: Teorías}
Explique un procedimiento que utilice resolución para determinar si una fórmula $\varphi$ cualquiera es tautología. Use su método para determinar si
$\varphi=(p\rightarrow q)\rightarrow(\neg p\rightarrow\neg q)$
es tautología.
\subsubsection*{Solución P2.}
Aquí va mi solución


\end{document}


